
We opened with an observation: in the Unterthiner CIFAR-10 CNN zoo, weight tensors predict
generalization gap better than they predict test accuracy --- a finding that inverts the conventional
assumption about which quantity is more legible from weights. Our experiments confirm and extend
this observation into three concrete results. Generalization gap is learnable from weight tensors
at Spearman r > 0.5 using standard encoders. Among the architectures we tested, NFT's cross-layer
attention achieves the highest gap Spearman (r = 0.575) with directional statistical advantage.
NFT's gap predictions contain information about true generalization gap that is independent of test
accuracy --- a partial Spearman of r = 0.73 --- revealing that weight tensors carry a distinct
overfitting signal that is not merely a shadow of absolute performance.

We investigated whether permutation-equivariant weight-space encoders show an advantage on
generalization gap prediction, and whether this advantage is specific to gap versus test accuracy.
Our controlled dual-target study on 10,000 CNNs produces four findings:

\begin{enumerate}
\item \textbf{Gap learnability established.} FlatMLP achieves Spearman r = 0.557 on generalization gap
\end{enumerate}
   prediction; DWSNet achieves r = 0.510. The A1 audit (Spearman(gap, $-test_acc) = -0.142)$
   confirms gap is not trivially derivable from test accuracy. Gap is a genuine, accessible
   weight-space signal.

\begin{enumerate}
\item \textbf{Cross-layer attention is the relevant inductive bias for gap.} NFT, the only tested encoder
\end{enumerate}
   with cross-layer attention, achieves the highest gap Spearman (r = 0.575, 95\% CI: [0.534, 0.616]).
   DWSNet (within-layer equivariance) and GNN (graph-structured equivariance) both underperform
   FlatMLP on gap --- equivariance alone is insufficient; the capacity to reason across layer
   boundaries provides advantage.

\begin{enumerate}
\item \textbf{NFT captures gap-specific overfitting structure.} The partial Spearman analysis
\end{enumerate}
   (r = 0.730, p = 1.$6\times 10^{-167})$ establishes that NFT extracts information about generalization gap
   that is statistically independent of test accuracy rank. Gap and test accuracy occupy distinct
   regions of the weight-space information landscape.

\begin{enumerate}
\item \textbf{Target-specific differential advantage not confirmed.} The $\Delta -based$ mechanism hypothesis is
\end{enumerate}
   not confirmed under our experimental conditions, confounded by an anomalously low FlatMLP
   test\_acc baseline (r = 0.279 vs. literature ~0.85). We identify the confounder, propose
   corrective experiments, and report the null result transparently.

\textbf{Future directions.} The most immediate priority is reproducing FlatMLP test\_acc Spearman $\approx$ 0.85
on the original Unterthiner zoo data with an extended 50-trial search budget, then recomputing $\Delta$
with a valid control baseline. The Schürholt PDFD zoo provides a multi-architecture evaluation venue
for replicating the NFT gap advantage and P3 partial Spearman. A within-layer-restricted NFT ablation
(h-m4, not executed) would directly test whether cross-layer attention is the operative mechanism
for gap advantage. A theoretical analysis connecting cross-layer weight co-variation to PAC-Bayes
gap magnitude could provide principled motivation for the empirical findings.

Weight tensors may be natural overfitting sensors --- encoding the distributed signature of
memorization across layer boundaries in a form that cross-layer attention architectures are well
positioned to read. The harder prediction problem is the one where weight space's structural
information proves most accessible.

